\documentclass{template}

\usepackage{mathtools}
\usepackage{stmaryrd}
\usepackage{bbm}

\usepackage{calrsfs}
\DeclareMathAlphabet{\pazocal}{OMS}{zplm}{m}{n}

\begin{document}

\section*{Master's Degree Proposals}

\subsection*{Lattice-based Folding Schemes}

\textbf{Background.} Classical elliptic curve-based cryptography enabled
numerous protocols and is arguably the cryptography branch that yields the most
efficient schemes (signature schemes such as Schnorr or ECDSA with sizes below
100~B, succinct proving systems based on elliptic curve pairing etc.) However,
such schemes are not secure against quantum computers. Consequently, the
question arises how to construct post-quantum secure counterparts of previous
schemes built on top of elliptic curves.

Currently lattice-based cryptography is the most promising solution to this
problem. Specifically, one of the major lattice-based assumptions (called
\textit{Short Integer Solution}) states that finding the ``short'' solution to
equation $\mathbf{Ax}=0$ for a random matrix $\mathbf{A} \in \mathbb{Z}_q^{n
\times m}$ with respect to $\mathbf{x} \in \mathbb{Z}_q^m$ is hard (where by
``short'' we mean $\|\mathbf{x}\|_{\infty}<\gamma$ for appropriate $\gamma$).
Many major tricks to making such schemes more efficient is related to
considering the cyclotomic number field $\mathbb{K} := \mathbb{Q}[X]/(X^n+1)$
for $n=2^k$, and using $R/qR$ instead of $\mathbb{Z}_q$
(where $R \cong \mathbb{Z}[X]/(X^n+1)$ is the ring of integers of
$\mathbb{K}$). That said, the equation $\mathbf{Ax}=0$ is now solved for a
random matrix in $R_q^{n \times m}$ and the norm is now taken with
respect to the geometry inherited from the structure of $\mathbb{K}$
(such modification is called \textit{Module Short Integer Solution} or MSIS for short).

\textbf{Goal.} There are currently numerous studies on how to apply such
assumptions on ``classical'' cryptographical constructions (such as digital
signatures or public key encryption). However, it is still unresolved how to put
more advanced protocols on lattice-based rails: multisignatures, threshold
signatures, succinct proving systems etc.

The promising line of research is \textit{folding schemes}. Recall that the goal
of proving systems over binary relation $\mathcal{R}$ is to demonstrate the
knowledge of the secret data (witness) $\mathbbm{w}$ such that $(\mathbbm{x},
\mathbbm{w}) \in \mathcal{R}$ for a public statement $\mathbbm{x}$. The
\textit{reduction} $\Pi: \mathcal{R} \to \mathcal{R}'$ is the protocol allowing
to reduce proving the statement $\mathbbm{x}$ over $\mathcal{R}$ to proving the
statement $\mathbbm{x}'$ over $\mathcal{R}'$ (with the property that any
efficient witness extractor of $\mathcal{R}'$ can extract the witness from
$\mathcal{R}$). Finally, the \textit{folding scheme} is a reduction
$\Pi_{\text{fd}}: \mathcal{R}_{\text{agg}} \times \mathcal{R}_{\text{comp}} \to
\mathcal{R}_{\text{agg}}$: i.e., instead of considering a complicated relation
$\mathcal{R}_{\text{agg}} \times \mathcal{R}_{\text{comp}}$ we reduce it to a
single relation $\mathcal{R}_{\text{agg}}$. This allows to fold proving 
thousands of same relations to a single relation $\mathcal{R}_{\text{agg}}$.

Turns out that lattice assumptions allow to build effective folding schemes.
Most recent papers on these topics are
\textit{LatticeFold}~\cite{cryptoeprint:2024/257} by Dan Boneh and Binyi Chen
(Stanford), \textit{Neo}~\cite{cryptoeprint:2025/294} by Wilson Nguyen and Srinath Setty
(Stanford and Microsoft Research), and \textit{Cyclo}~\cite{cryptoeprint:2026/359} by
Albert Garreta et al. (University of Tartu and Nethermind). 

\textbf{Goal \#1.} Survey all aforementioned schemes and understand underlying techniques.

\textbf{Goal \#2.} The schemes presented are \textit{not zero-knowledge}. 
The goal is to build the zero-knowledge-secure version of one of presented protocols
and prove it.

\textit{Details.} Two primary reasons why they are not zero-knowledge:
\begin{itemize}
    \item Underlying sum-check protocol (\textit{i.e.}, showing
    $\sum_{\vec{\mathbf{b}} \in \{0,1\}^{\nu}} f(\vec{\mathbf{b}}) = H$ for $f \in R[X_1,\dots,X_{\nu}]$ and public $H \in R$ for some ring $R$) is not
    zero-knowledge by default. One should apply techniques to make it
    zero-knowledge (\textit{e.g.}, apply techniques from
    \cite{cryptoeprint:2017/305}).
    \item Ajtai commitment scheme $\mathbf{x} \mapsto \mathbf{Ax}$ is
    \textit{binding} under MSIS assumption, but not hiding. Is it possible to
    use randomized (hiding) Ajtai commitment scheme $(\mathbf{x},\mathbf{r}) \mapsto
    \mathbf{Ax}+\mathbf{Br}$ instead? How much worse schemes become in terms of efficiency?
\end{itemize}

\textbf{Goal \#3.} One concrete practical application of folding scheme is
batching multiple \textit{range proofs} (proofs that some private amount $s \in
\mathbb{N}$ lies in range $0 \leq s \leq L$ for $L \in \mathbb{N}$). Is it possible
to specifically optimize the scheme where $\mathcal{R}_{\text{comp}}$ is the range proof?

\textit{Possible solution.} Apply techniques from \cite{cryptoeprint:2020/517} and see what happens.

\subsection*{Isogeny-based Cryptography}

Another quite novel area of post-quantum research requiring
arguably the most heavy algebraic preliminaries is \textit{isogeny-based cryptography}.

\textbf{Goal \#1.} Understand the SQISign construction:

\begin{center}
    \url{https://sqisign.org/spec/sqisign-20250707.pdf}
\end{center}

\textbf{Goal \#2.} Is it possible to construct key derivation, multisignatures, threshold signatures on top?

\bibliographystyle{alphaurl}
\bibliography{refs-proposal}

\end{document}